\section{Rising-continuity, frames, and products of F-series}
\label{sec:rising-fseries}

The 2022 dissertation, the 2023 paper on variable topologies, and the 2025
paper on products all study series whose target topology depends on the input
digit sequence.  The stable core is: canonical truncation limits,
point-selected completions (frames), digit-multiplicative F-series, and
finite-level Fourier recursions.  Several Banach-algebra and product-of-
distributions conclusions in the sources require repairs recorded below.

\subsection{Rising-continuity}

Let $p$ be prime, let $K$ be a complete nonarchimedean field, and let
$f:\Z_p\to K$.

\begin{definition}[rising-continuity]
The function $f$ is \emph{rising-continuous} when
\begin{equation}
\label{eq:rising-continuity}
 f(z)=\lim_{N\to\infty}f([z]_{p^N})
 \qquad(z\in\Z_p).
\end{equation}
Write $\widetilde C(\Z_p,K)$ for the class of such functions.
\end{definition}

For an arbitrary function $f:\Z_p\to K$, define the source's level
differences by
\begin{align}
 (\nabla_{p^0}f)(z)&=f(0),\label{eq:del-zero}\\
 (\nabla_{p^N}f)(z)
 &=f([z]_{p^N})-f([z]_{p^{N-1}}),
 &&N\geq1,\label{eq:del-positive}
\end{align}
and its extended del quantity by
\begin{equation}
\label{eq:del-norm}
 \|f\|_{\nabla}
 =\sup_{N\geq0}\sup_{z\in\Z_p}|(\nabla_{p^N}f)(z)|_K
 \in[0,\infty].
\end{equation}

\begin{proposition}[the exact domain of the del norm]
\label{prop:del-norm-domain}
If $f\in\widetilde C(\Z_p,K)$, then
\begin{equation}
\label{eq:del-equals-sup}
 \|f\|_{\nabla}=\sup_{z\in\Z_p}|f(z)|_K
\end{equation}
as an equality in $[0,\infty]$.  Consequently
\[
 \widetilde C_b(\Z_p,K)
 :=\{f\in\widetilde C(\Z_p,K):\|f\|_\nabla<\infty\}
\]
is precisely the bounded rising-continuous class, and it is a Banach space
under $\|\cdot\|_\nabla$.
\end{proposition}

\begin{proof}
For every $N$ and $z$, the ultrametric inequality bounds
$|(\nabla_{p^N}f)(z)|_K$ by $\sup_w|f(w)|_K$, including $N=0$.  Conversely,
the finite telescoping identity
\[
 f([z]_{p^N})=\sum_{h=0}^N(\nabla_{p^h}f)(z)
\]
and the ultrametric inequality give
$|f([z]_{p^N})|_K\leq\|f\|_\nabla$.  Rising-continuity and continuity of
the absolute value let $N\to\infty$, proving the reverse inequality in
\cref{eq:del-equals-sup}.  Thus the two finite norms agree.  Bounded
$K$-valued functions are complete under the supremum norm.  If $(f_m)$ is a
uniformly Cauchy sequence of bounded rising-continuous functions with uniform
limit $f$, then for fixed $z$ the ultrametric estimate
\[
 |f(z)-f([z]_{p^N})|_K
 \leq\max\{|f(z)-f_m(z)|_K,
            |f_m(z)-f_m([z]_{p^N})|_K,
            |f_m([z]_{p^N})-f([z]_{p^N})|_K\}
\]
first becomes small by choosing $m$ and then $N$.  Hence $f$ is
rising-continuous, so the displayed subspace is closed and complete.
\end{proof}

Every continuous function is rising-continuous because
$[z]_{p^N}\to z$ in $\Z_p$.  The converse is false: the condition tests one
distinguished convergent sequence at each point, not every convergent
sequence.  For every odd prime $q$, the shortened numen
$\chi_q:\Z_2\to\Z_q$ is the principal example (the preceding definition with
$p=2$ and $K=\Q_q$): it is rising-continuous everywhere but discontinuous at
each nonnegative integer.

Define the truncation of $f$ at level $N$ by
\begin{equation}
\label{eq:function-truncation}
 f_N(z)=f([z]_{p^N})
 =\sum_{k=0}^{p^N-1}f(k)\ind{z\equiv k\pmod{p^N}}.
\end{equation}

\begin{proposition}[pointwise versus uniform truncation]
\label{prop:rising-vs-continuous}
For an arbitrary $f:\Z_p\to K$:
\begin{enumerate}[label=\textup{(\roman*)}]
\item $f_N\to f$ pointwise if and only if $f$ is rising-continuous;
\item $f_N\to f$ uniformly if and only if $f$ is continuous.
\end{enumerate}
\end{proposition}

\begin{proof}
Part (i) is the definition.  Every $f_N$ is locally constant and hence
continuous, so a uniform limit is continuous.  Conversely, a continuous
function on compact $\Z_p$ is uniformly continuous.  Since
$|z-[z]_{p^N}|_p\leq p^{-N}$ uniformly in $z$, uniform continuity gives
$\sup_z|f(z)-f_N(z)|\to0$.
\end{proof}

\subsection{Van der Put paths}

Set $\lambda_p(0)=0$ and
\[
 \lambda_p(n)=1+\lfloor\log_p n\rfloor\qquad(n\geq1).
\]
For $n\geq1$, let $n^-$ be the integer obtained by deleting the highest
nonzero base-$p$ digit of $n$.  Define
\begin{align}
 v_0(z)&=1,\\
 v_n(z)&=\ind{z\equiv n\pmod{p^{\lambda_p(n)}}},\qquad n\geq1,\label{eq:vdp-basis}\\
 c_0(f)&=f(0),\\
 c_n(f)&=f(n)-f(n^-),\qquad n\geq1.
\end{align}

\begin{theorem}[rising van der Put representation]
\label{thm:rising-vdp}
A function $f:\Z_p\to K$ is rising-continuous if and only if
\begin{equation}
\label{eq:rising-vdp}
 f(z)=\sum_{n=0}^{\infty}c_n(f)v_n(z)
\end{equation}
converges at every $z$.  Equivalently, along every truncation path,
\begin{equation}
\label{eq:vdp-path-condition}
 c_{[z]_{p^N}}(f)\,ind{\lambda_p([z]_{p^N})=N}
 \longrightarrow0.
\end{equation}
The expansion is unique.
\end{theorem}

\begin{proof}
At a fixed $z$, exactly one possible new van der Put ball appears when the
$N$th digit is nonzero: its index is $[z]_{p^N}$ and its coefficient is the
difference between two successive effective truncation values.  Therefore
the partial sum through level $N$ telescopes to $f([z]_{p^N})$.  In a
complete nonarchimedean field, the path series converges exactly when its
terms tend to zero, giving \cref{eq:vdp-path-condition}; its sum is $f(z)$
exactly when \cref{eq:rising-continuity} holds.  Evaluation successively on
the nested balls recovers every coefficient, proving uniqueness.
\end{proof}

\begin{corollary}[unique rising continuation]
\label{cor:unique-rising-continuation}
Let $g:\N_0\to K$.  A rising-continuous extension to $\Z_p$ exists if and
only if $g([z]_{p^N})$ converges in $K$ for every $z\in\Z_p$.  When it exists
it is unique and equals
\[
 \widetilde g(z)=\lim_{N\to\infty}g([z]_{p^N}).
\]
\end{corollary}

\begin{proof}
If a rising-continuous extension $f$ exists, then
$g([z]_{p^N})=f([z]_{p^N})\to f(z)$ for every $z$, so the displayed limit is
necessary and determines $f$ uniquely.  Conversely, suppose every displayed
limit exists and use it to define $\widetilde g$.  If $m\in\N_0$, then
$[m]_{p^N}=m$ for all sufficiently large $N$, hence
$\widetilde g(m)=g(m)$.  Therefore
\[
 \widetilde g([z]_{p^N})=g([z]_{p^N})\longrightarrow\widetilde g(z),
\]
which is exactly rising-continuity.  Thus $\widetilde g$ is the required
extension, and the necessity argument already proved uniqueness.
\end{proof}

This is the precise uniqueness condition missing from several functional-
equation statements in the 2026 manual: uniqueness on $\N_0$ does not by
itself imply uniqueness on all infinite digit strings, but rising-continuity
does.

\subsection{Frames: a topology selected at each point}

Let $F$ be a global field and let $\operatorname{Comp}(F)$ denote a chosen
collection of complete valued extensions of $F$.

\begin{definition}[frame]
An $F$-\emph{frame} on a set $X$ is a map
\[
 \mathcal F:X\to\operatorname{Comp}(F).
\]
A function $f$ is compatible with $\mathcal F$ when
$f(x)\in\mathcal F(x)$ for every $x$.  A sequence $f_N$ converges to $f$ in
the frame when
\begin{equation}
\label{eq:frame-convergence}
 f_N(x)\longrightarrow f(x)\quad\text{in }\mathcal F(x)
 \qquad(x\in X).
\end{equation}
\end{definition}

Thus a frame is not one exotic topology on a disjoint union.  It is a rule
for which ordinary completion is to judge convergence at each input.
Pointwise sums and products of compatible functions are meaningful because
both operands at a fixed point lie in the same selected field.

The standard two-place example in the body of the 2023 paper uses complex
convergence on $\N_0$ and $q$-adic convergence on
$\Z_p\setminus\N_0$.  Its abstract describes the assignments in the opposite
order; the body definition controls this edition.

\begin{sourceissue}[an unrestricted frame function need not have a sup norm]
One source proposition declares all functions from compact $X$ to a complete
field to be a Banach space under the supremum norm.  Compactness bounds
continuous functions, not arbitrary functions.  The statement becomes
valid after restricting to bounded functions (and proving completeness
pointwise plus uniformly), or to continuous functions in a fixed target.
The same issue reappears in the dissertation's proposed Banach algebra of all
rising-continuous functions: rising-continuity alone does not imply
boundedness or finiteness of the displayed ``del'' norm.
\end{sourceissue}

\subsection{Digit-multiplicative F-series}

For a digit prefix of fixed length $N$, use the padded count
\begin{equation}
\label{eq:padded-count}
 \#^{(N)}_{p:j}(z)=\#\{0\leq h<N:d_h(z)=j\}.
\end{equation}
This count includes high zero digits.  The ordinary digit count of the
integer $[z]_{p^N}$ silently discards them and is not interchangeable with
\cref{eq:padded-count} when a zero-digit multiplier occurs.

Fix a complete valued field $E$ containing nonzero parameters
$d,q_0,\ldots,q_{p-1}$.  For $z\in\Z_p$, put
\[
 u_N(z)=d^{-N}\prod_{j=0}^{p-1}
 q_j^{\#^{(N)}_{p:j}(z)},
 \qquad
 D_E=\left\{z\in\Z_p:\sum_{N\geq0}u_N(z)
                       \text{ converges in }E\right\}.
\]
Define the $E$-valued F-series on its exact convergence domain by
\begin{equation}
\label{eq:general-F-series}
 S^E_{d;\mathbf q}:D_E\longrightarrow E,
 \qquad S^E_{d;\mathbf q}(z)=\sum_{N=0}^{\infty}u_N(z).
\end{equation}
 For $0\leq k<p$, the digit identity
\[
 \#^{(N+1)}_{p:j}(pz+k)
 =\#^{(N)}_{p:j}(z)+\ind{j=k}
\]
gives the termwise tail identity
\[
 u_{N+1}(pz+k)=\frac{q_k}{d}u_N(z).
\]
Since $q_k/d\ne0$, it follows without an unstated convergence condition that
$pz+k\in D_E$ if and only if $z\in D_E$, and on this common domain
\begin{equation}
\label{eq:F-series-FE}
 S^E_{d;\mathbf q}(pz+k)=1+\frac{q_k}{d}S^E_{d;\mathbf q}(z).
\end{equation}
A frame may select a different completion at different inputs, but an
identity comparing $z$ with $pz+k$ is typed only after both values are placed
in one specified field $E$ (or after a specified identification of their
target fields).

The primitive numen series \cref{eq:partial-F-series} is the affine version:
its $N$th term contains a digit-dependent translation $c_{d_N}$ and the
ordered product of all earlier linear multipliers.  Fix a complete valued
field $E$ and scalars $a_0,\ldots,a_{p-1},b_0,\ldots,b_{p-1}\in E$.  In this
scalar commuting case put
\[
 D_X=\left\{z\in\Z_p:
 \sum_{N=0}^{\infty}b_{d_N(z)}
  \prod_{h=0}^{N-1}a_{d_h(z)}\text{ converges in }E\right\}
\]
and define
\begin{equation}
\label{eq:scalar-affine-F-series}
 X:D_X\longrightarrow E,\qquad
 X(z)=\sum_{N=0}^{\infty}b_{d_N(z)}
 \prod_{j=0}^{p-1}a_j^{\#^{(N)}_{p:j}(z)}.
\end{equation}
For $0\leq k<p$ and $z\in D_X$ one has $pz+k\in D_X$ and
\begin{equation}
\label{eq:scalar-affine-FE}
 X(pz+k)=a_kX(z)+b_k.
\end{equation}
Conversely, if $a_k\ne0$, then $pz+k\in D_X$ implies $z\in D_X$.
Indeed, the zeroth summand at $pz+k$ is $b_k$, and its summand of index
$N+1$ is $a_k$ times the summand of index $N$ at $z$.  Thus convergence in
the forward direction is automatic even when $a_k=0$, while division by
$a_k$ in the converse is valid exactly under the displayed nonvanishing
hypothesis.

For the Collatz numen, the 2023 paper writes the useful frame identity
\begin{equation}
\label{eq:chi3-F-series}
 \chi_3(z)=-\frac12+\frac14
 \sum_{N=0}^{\infty}\frac{3^{\#^{(N)}_{2:1}(z)}}{2^N}.
\end{equation}
On nonnegative integers the right side is read archimedeanly; on
$\Z_2\setminus\N_0$ it is read $3$-adically.  This differing-codomain
identity does not follow from \cref{cor:unique-rising-continuation}, whose
domain and codomain are fixed.  It follows directly from one finite identity.
Indeed, write
\[
 A_N(z)=\frac{3^{\#^{(N)}_{2:1}(z)}}{2^N},
 \qquad S_N(z)=\sum_{h=0}^{N-1}A_h(z),
\]
and let $d_h(z)\in\{0,1\}$ be the $h$th digit.  The $N$-digit numen partial
sum is
\[
 \chi_{3,N}(z)=\frac12\sum_{h=0}^{N-1}d_h(z)A_h(z).
\]
Since $A_{h+1}=(\frac12+d_h)A_h$ and $A_0=1$, telescoping gives the exact
finite relation
\begin{equation}
\label{eq:chi3-finite-frame-identity}
 \chi_{3,N}(z)=\frac14S_N(z)-\frac12+\frac12A_N(z).
\end{equation}
If $z\in\N_0$, its one-count is eventually constant and $A_N(z)\to0$ in
$\R$.  If $z\in\Z_2\setminus\N_0$, its binary expansion has infinitely many
one digits, so $|A_N(z)|_3=3^{-\#^{(N)}_{2:1}(z)}\to0$.  In the first case
$\chi_{3,N}(z)$ is eventually the finite numen value; in the second it
converges in $\Q_3$ by the primitive numen series.  Passing to the limit in
\cref{eq:chi3-finite-frame-identity} proves
\cref{eq:chi3-F-series} separately in the stated completion at every input.

\begin{sourceissue}[pointwise convergence claim in the 2023 paper]
One necessity proof chooses a general nonnatural input and then replaces it
by $z=-1$.  That proves a global obstruction, not the asserted pointwise
necessity for the original digit sequence.  Placewise convergence should be
checked from the actual prefix products, as in
\cref{prop:numen-pointwise-convergence}, rather than from that substitution.
\end{sourceissue}

\subsection{Products and the exact finite recursion}

Let $R$ be a commutative ring with identity, let
$a_{j,k},b_{j,k}\in R$, and let
$X_1,\ldots,X_d:\Z_p\to R$ be degree-one digit-recursive functions
satisfying
\begin{equation}
\label{eq:degree-one-recursion}
 X_j(pz+k)=a_{j,k}X_j(z)+b_{j,k}
 \qquad(1\leq j\leq d,\ 0\leq k<p).
\end{equation}
For a multi-index $\mathbf n=(n_1,\ldots,n_d)\in\N_0^d$, put
\[
 X_{\mathbf n}=\prod_{j=1}^dX_j^{n_j}.
\]
If $\mathbf m\leq\mathbf n$ coordinatewise, define
\begin{align}
 \binom{\mathbf n}{\mathbf m}&=\prod_j\binom{n_j}{m_j},\\
 a_{\mathbf m,k}&=\prod_ja_{j,k}^{m_j},\\
 b_{\mathbf n-\mathbf m,k}&=\prod_jb_{j,k}^{n_j-m_j},\\
 r_{\mathbf m,\mathbf n,k}
 &=\binom{\mathbf n}{\mathbf m}
   a_{\mathbf m,k}b_{\mathbf n-\mathbf m,k}.
\end{align}

\begin{proposition}[multinomial branch recursion]
\label{prop:product-branch-recursion}
For every $k$,
\begin{equation}
\label{eq:product-branch-recursion}
 X_{\mathbf n}(pz+k)
 =\sum_{\mathbf m\leq\mathbf n}
 r_{\mathbf m,\mathbf n,k}X_{\mathbf m}(z).
\end{equation}
\end{proposition}

\begin{proof}
Raise each equation in \cref{eq:degree-one-recursion} to the power $n_j$,
expand by the binomial theorem, multiply over $j$, and collect the term with
multi-index $\mathbf m$.
\end{proof}

This finite algebraic identity is independent of every limiting assumption
in the 2025 paper.

\subsection{Finite Fourier recursion and formal solutions}

Fix through the end of \cref{prop:finite-transform-limit} a complete
nonarchimedean field $K$ of residue
characteristic $q\ne p$, a compatible system of $p$-power roots of unity,
and the normalized $K$-valued Haar integral of
\cref{sec:pq-fourier}.  Assume the branch parameters and the functions under
transformation are $K$-valued.  A point-varying frame by itself is not a
single coefficient field in which Fourier sums can be added.

Let
\begin{equation}
\label{eq:alpha-product}
 \alpha_{\mathbf m,\mathbf n}(t)
 =\frac1p\sum_{k=0}^{p-1}
 r_{\mathbf m,\mathbf n,k}\langle-t,k\rangle,
 \qquad
 \alpha_{\mathbf n}=\alpha_{\mathbf n,\mathbf n}.
\end{equation}
Fix $\mathbf n$ and assume explicitly that
$X_{\mathbf m}\in C(\Z_p,K)$ for every $\mathbf m\leq\mathbf n$.  Then all
Fourier transforms below are the normalized Haar integrals of
\cref{eq:pq-Fourier}.  The substitution
$x=pz+k:\Z_p\xrightarrow{\sim}k+p\Z_p$ gives
\[
 \int_{k+p\Z_p}f(x)\,dx
 =\frac1p\int_{\Z_p}f(pz+k)\,dz
\]
for continuous $f$.  Applying this identity to
$f(x)=X_{\mathbf n}(x)\langle-t,x\rangle$, using
\cref{eq:product-branch-recursion}, and interchanging only finite sums gives
\begin{equation}
\label{eq:product-Fourier-recursion}
 \widehat X_{\mathbf n}(t)
 =\sum_{\mathbf m\leq\mathbf n}
 \alpha_{\mathbf m,\mathbf n}(t)
 \widehat X_{\mathbf m}(pt).
\end{equation}
At finite truncation level the identity is exact before any limiting
integrability claim, but it has a necessary one-level shift.  Put
\[
 X_{j,N}(z)=X_j([z]_{p^N}),\qquad
 X_{\mathbf n,N}(z)=\prod_{j=1}^dX_{j,N}(z)^{n_j}.
\]
Then
\begin{equation}
\label{eq:finite-product-Fourier-recursion}
 \widehat X_{\mathbf n,N+1}(t)
 =\sum_{\mathbf m\leq\mathbf n}
 \alpha_{\mathbf m,\mathbf n}(t)\widehat X_{\mathbf m,N}(pt).
\end{equation}
Indeed, $[pz+k]_{p^{N+1}}=p[z]_{p^N}+k$.  More generally, for
$f_N(z)=f([z]_{p^N})$,
\begin{equation}
\label{eq:finite-truncation-transform}
 \widehat f_N(t)=
 \frac{\ind{p^Nt=0}}{p^N}
 \sum_{r=0}^{p^N-1}f(r)\langle-t,r\rangle.
\end{equation}

At $t=0$, separate the diagonal term in
\cref{eq:product-Fourier-recursion}:
\begin{equation}
\label{eq:formal-mean-equation}
 (1-\alpha_{\mathbf n}(0))\widehat X_{\mathbf n}(0)
 =\sum_{\mathbf m<\mathbf n}
 \alpha_{\mathbf m,\mathbf n}(0)\widehat X_{\mathbf m}(0).
\end{equation}

\begin{proposition}[formal compatibility classification]
\label{prop:formal-compatibility}
Assume the lower-degree transforms are fixed.
\begin{enumerate}[label=\textup{(\roman*)}]
\item If $\alpha_{\mathbf n}(0)\ne1$, the formal recursion has a unique
      value at $0$, and hence a unique value at every $t\in\Gamma_p$ by
      repeated use of \cref{eq:product-Fourier-recursion} until $p^Nt=0$.
\item If $\alpha_{\mathbf n}(0)=1$ and the right side of
      \cref{eq:formal-mean-equation} is nonzero, no formal solution exists.
\item If $\alpha_{\mathbf n}(0)=1$ and that right side is zero, the value
      $\widehat X_{\mathbf n}(0)$ is free and the solutions form a
      one-dimensional affine family.
\end{enumerate}
\end{proposition}

\begin{proof}
\Cref{eq:formal-mean-equation} gives the three alternatives at
$0$.  Every $t\in\Gamma_p$ has $p$-power order, so $p^Nt=0$ for some $N$;
running \cref{eq:product-Fourier-recursion} backward from that terminal value
determines all other values.
\end{proof}

The paper calls the locus $\alpha_{\mathbf n}(0)=1$ a breakdown locus.
\Cref{prop:formal-compatibility} shows why breakdown does not automatically
mean nonexistence: the compatibility sum decides between no solution and an
affine family.

\subsection{M-functions and limits}

For scalar branch multipliers $r_0,\ldots,r_{p-1}\in K^\times$, the associated
M-function is the prefix cocycle
\begin{equation}
\label{eq:M-function}
 M_N(z)=\prod_{h=0}^{N-1}r_{d_h(z)}
 =r_0^N\prod_{j=1}^{p-1}\left(\frac{r_j}{r_0}\right)^{
 \#^{(N)}_{p:j}(z)}.
\end{equation}
It satisfies
\[
 M_{N+L}(z)=M_N(z)M_L(\sh^{\circ N}z).
\]
These cocycles are the proposed majorants for passing from
\cref{eq:finite-truncation-transform} to a framewise distributional limit.

Before stating the exact limit consequence, set
\begin{align*}
 \mathcal S(\Z_p,K)
 &=\{\phi:\Z_p\to K:\phi\text{ is locally constant modulo }p^N
       \text{ for some }N\},\\
 \mathcal S(\Z_p,K)'
 &=\operatorname{Hom}_K(\mathcal S(\Z_p,K),K).
\end{align*}

\begin{proposition}[the exact consequence of transform convergence]
\label{prop:finite-transform-limit}
Fix $\mathbf n$.  Suppose that for every $\mathbf m\leq\mathbf n$ and every
$t\in\Gamma_p$ the limit
\[
 A_{\mathbf m}(t)=\lim_{N\to\infty}
 \widehat X_{\mathbf m,N}(t)
\]
exists in the fixed complete field $K$.  Then
\begin{equation}
\label{eq:limit-product-Fourier-recursion}
 A_{\mathbf n}(t)=
 \sum_{\mathbf m\leq\mathbf n}
 \alpha_{\mathbf m,\mathbf n}(t)A_{\mathbf m}(pt)
 \qquad(t\in\Gamma_p).
\end{equation}
Moreover, every function $A:\Gamma_p\to K$ defines a unique algebraic
distribution $T_A\in\mathcal S(\Z_p,K)'$ by
\begin{equation}
\label{eq:formal-Fourier-distribution}
 \langle T_A,\phi\rangle
 =\sum_{t\in\Gamma_p}\widehat\phi(t)A(-t),
 \qquad \phi\in\mathcal S(\Z_p,K),
\end{equation}
where the sum is finite.  With the Fourier sign of
\cref{eq:pq-Fourier}, this convention gives
$\langle T_A,\langle-t,\cdot\rangle\rangle=A(t)$.  In particular, if
\[
 \langle\mu_{\mathbf m,N},\phi\rangle
 =\int_{\Z_p}X_{\mathbf m,N}(z)\phi(z)\,dz,
\]
then $\mu_{\mathbf m,N}\to T_{A_{\mathbf m}}$ pointwise on
$\mathcal S(\Z_p,K)$ for every $\mathbf m\leq\mathbf n$.
\end{proposition}

\begin{proof}
Take $N\to\infty$ in
\cref{eq:finite-product-Fourier-recursion}.  Its sum over
$\mathbf m\leq\mathbf n$ is finite, so convergence of the displayed terms is
all that is required and yields
\cref{eq:limit-product-Fourier-recursion}.  Every locally constant test
function has a unique finite character expansion
$\phi(z)=\sum_t\widehat\phi(t)\langle t,z\rangle$ by finite Fourier
inversion.  Hence \cref{eq:formal-Fourier-distribution} is a well-defined
$K$-linear functional, and evaluating it on the character
$\langle-t,\cdot\rangle$ proves the final identity and uniqueness.  Finally,
finite Fourier inversion for $\phi$ gives
\[
 \langle\mu_{\mathbf m,N},\phi\rangle
 =\sum_t\widehat\phi(t)\widehat X_{\mathbf m,N}(-t),
\]
where the sum is finite.  Termwise passage to the limit is therefore exact
and yields \cref{eq:formal-Fourier-distribution} with
$A=A_{\mathbf m}$.
\end{proof}

\begin{status}[conditional analytic theorem in the 2025 paper]
The paper explicitly introduces a Main Limit Assumption before asserting
that the formal Fourier data are realized by frame-quasi-integrable products.
\Cref{prop:finite-transform-limit} isolates exactly what follows if all of
the required finite-transform limits exist in one fixed completion.

The later Main Limit Lemma proposes sufficient conditions in terms of strong
frame summability of $M_N$, uniform control of shifted $X$, and
$(\mathcal F,M)$-type decay.  The printed proof does not establish those
sufficient conditions, so the ensuing base case, induction, power theorem,
and main product corollary remain conditional.
\end{status}

\subsection{Test functions, distributions, and what the product means}

More generally, for a commutative coefficient ring $R$ with identity, let
\begin{align}
 \mathcal S(\Z_p,R)
 &=\{\phi:\Z_p\to R:\phi\text{ is locally constant modulo }p^N
       \text{ for some }N\},\\
 \mathcal S(\Z_p,R)'&=\operatorname{Hom}_R(\mathcal S(\Z_p,R),R).
\end{align}
The second module is the algebraic space of $(p,R)$-adic distributions.  If
$\mu$ is a distribution and $\phi$ a test function, multiplication by a test
function is unambiguous:
\[
 \langle\phi\mu,\psi\rangle
 =\langle\mu,\phi\psi\rangle
 \qquad\bigl(\psi\in\mathcal S(\Z_p,R)\bigr).
\]
Here the distribution occupies the first slot and the test function the
second.  This convention is used throughout the subsection.
Products of two arbitrary distributions are not defined by this rule.

\begin{sourceissue}[the product algebra is transported, not intrinsic]
The 2025 paper defines, for its designated family,
\[
 (X^m(z)\,dz)(X^n(z)\,dz):=X^{m+n}(z)\,dz
\]
and defines the corresponding ``convolution powers'' to be the recursively
constructed formal Fourier data.  The exact formal object is the polynomial
algebra $R[T]$, with $T^n$ labelled by the symbol $[X^n(z)\,dz]$ and with
$T^mT^n=T^{m+n}$.  This construction exists without realizing any symbol as
a distribution.

Suppose separately that distributions
$\mu_n\in\mathcal S(\Z_p,R)'$ have actually been constructed for the named
symbols, and define the $R$-linear realization map
\[
 \Phi:R[T]\longrightarrow\mathcal S(\Z_p,R)',
 \qquad \Phi\!\left(\sum_na_nT^n\right)=\sum_na_n\mu_n.
\]
The rule
\begin{equation}
\label{eq:transported-product-rule}
 \mu_m\mathbin\odot\mu_n:=\mu_{m+n}
\end{equation}
descends to a well-defined multiplication on $\operatorname{im}\Phi$ if and
only if $\ker\Phi$ is an ideal of $R[T]$.  Indeed, well-defined descent is
equivalent to
$\Phi(f)=\Phi(f')\Rightarrow\Phi(fg)=\Phi(f'g)$ for every $g$, which is
equivalent to closure of $\ker\Phi$ under multiplication by every
polynomial.  When this holds, the realized algebra is canonically
$R[T]/\ker\Phi$; when $\Phi$ is injective it is the polynomial algebra
itself.  If $\ker\Phi$ is not an ideal, the declaration
\cref{eq:transported-product-rule} depends on the chosen polynomial
representative and is not a product on the image.

The paper also warns that ordinary convolution of its Fourier data generally
diverges.  The formal product above is therefore not a proof that arbitrary
distributions possess pointwise products, nor that standard Fourier
convolution realizes the declared multiplication.
\end{sourceissue}

\subsection{Failure and repair ledger for the product paper}

The following defects affect theorem status rather than notation alone.
\begin{enumerate}[label=\textup{(\arabic*)}]
\item Haar probability of a residue class modulo $p^N$ is printed in two
      places as $1/p$; it is $p^{-N}$.  The corresponding finite character
      expansion needs the same correction.
\item In the archimedean Wiener algebra, absolutely summable Fourier
      coefficients must carry the $\ell^1$ norm
      $\sum_t|\widehat f(t)|$, not the coefficient supremum norm printed in
      the paper.
\item The Main Limit Lemma repeatedly infers that a positive sequence tending
      to zero is monotone on a tail, and that a tail sum is
      $O$ of its first term.  Neither follows from convergence or from a root
      bound.  The base case also uses the false estimate $Na_N=O(a_N)$.
\item The localization $R[(1-a_0)^{-1}]$ can be quotiented by an ideal $I$
      only after requiring $I$ to be disjoint from every power of
      $1-a_0$; the weaker printed condition $1-a_0\notin I$ does not suffice.
      A fraction field of the quotient additionally requires a domain, for
      example a prime ideal.  Multivariable polynomial localizations are not
      Dedekind domains in dimension greater than one.
\item A final corollary says a breakdown index gives no solution for any
      initial data, contradicting the compatible affine-family case already
      proved in the paper and restated in
      \cref{prop:formal-compatibility}.
\end{enumerate}

A repair programme must first correct Haar and Wiener normalization, impose
domain-safe localization hypotheses, replace the false monotonicity steps by
a genuine monotone majorant or regular-variation estimate, and prove rising
continuity of the proposed limits.  Only then can the analytic product and
inversion conclusions be promoted from conditional claims.
